\begin{table}[tbp]
    \centering
    \footnotesize
    \setlength{\tabcolsep}{3pt}
    \caption{Dinâmica de depleção vesicular e correlação temporal de disparo no Experimento 1.}
    \label{tab:resultados_exp1_dinamica}
    \begin{tabular}{lcccc}
        \toprule
        \textbf{Padrão} & \textbf{$R_{\text{ss}}$} & \textbf{$R_{\text{mín}}$} & \textbf{Correlação Pop. ($r$)} & \textbf{Correlação Célula ($r$)} \\
        \midrule
        Padrão 1 ($S_1$) & $0{,}1058$ & $0{,}0605$ & \multirow{2}{*}{$r = -0{,}3035$} & \multirow{2}{*}{$-0{,}2705 \pm 0{,}1794$} \\
        Padrão 2 ($S_2$) & $0{,}1223$ & $0{,}0607$ & & \\
        \midrule
        \textbf{Regime} & \multicolumn{2}{c}{Depleção estacionária ($\approx 89\%$)} & \multicolumn{2}{c}{\textbf{Defasagem temporal estrita ($r < 0$)}} \\
        \bottomrule
    \end{tabular}
    \vspace{0.2em}
    \begin{minipage}{0.95\textwidth}
        \footnotesize \textbf{Nota:} $R_{\text{ss}}$ e $R_{\text{mín}}$ representam a fração vesicular estacionária média e mínima de Tsodyks-Markram ($\tau_{\text{rec}} = \qty{600}{\milli\second}$, $U = 0{,}45$). A correlação linear de Pearson ($r$) foi calculada sobre os perfis temporais médios de disparo (PSTH com 10 bins de $\qty{14}{\milli\second}$ do ciclo teta de $\qty{140}{\milli\second}$), comprovando que os disparos ocorrem em fases descorrelacionadas entre $S_1$ e $S_2$.
    \end{minipage}
\end{table}
